% Introduction: VVUQ primer, the single-humanoid thesis, the VVUQ-01 lineage,
% scope and the regulatory moment, contributions and roadmap. Prose only; file
% names are deferred to the Methods tables, and the literature is cited rather
% than described.

Verification, validation, and uncertainty quantification, abbreviated VVUQ, is
the discipline that decides whether a computational result is credible enough to
act on. In the plain language of an oncology clinical trial the three questions
are concrete. Verification asks whether the code was built correctly, that the
software computes what its specification says it computes. Validation asks
whether the right model was built, that the computed behavior agrees with an
independent reference of the real world. Uncertainty quantification asks how
confident one should be in the answer, given the spread of results across
repeated runs and the assumptions that produced them. These are the same three
questions a trial statistician asks of an endpoint, a sponsor asks of a safety
signal, and a regulator asks of any model offered as evidence. The survey
vocabulary for digital twins in precision medicine frames VVUQ in exactly this
way \cite{sel2025vvuq}, and the discipline already has direct oncology precedent
in the verification, validation, and uncertainty quantification of
physiologically based pharmacokinetic models used for treatment planning
\cite{zaid2025pbpk}. What distinguishes the present work is that each of the
three questions is not an abstraction but an executed check: a verification
predicate that must pass, a validation comparison against a committed reference,
and a dispersion bound computed across seeded runs, composed into a single gate
that returns accept, block, or escalate. The reader should hold that picture
throughout, because the entire argument of this paper rests on the claim that
running those checks well, not generating the code they judge, is where the work
and the risk actually live.

That claim is the thesis. For a single autonomous surgical humanoid the robotic
code assurance process, not code generation, is the substantial and
decision-bearing part of the artificial intelligence workflow, and holding VVUQ
to a higher standard than the code itself is what will make upcoming physical AI
oncology trial deliverables faster, less expensive, and more beneficial to
patients than conventional verifications. The asymmetry is sharpened by the
embodiment. A teleoperated multi-arm cart distributes its action across several
instruments under a surgeon's continuous control, with a bedside assistant and a
second cart available to take over. A single humanoid concentrates all surgical
action into two hands, with no second cart and no human bedside assistant in the
loop, so all error potential is concentrated in one body. The consequence is
structural rather than rhetorical: the assurance layer must be larger, stricter,
and more thoroughly exercised than the control code it judges, and any candidate
behavior must clear ten distinct humanoid-specific gates before it may ship. Code
generation, by contrast, is fast and deterministic; a behavior model or a
structured results file is produced in microseconds and reproduced exactly from a
fixed seed. The decision to ship is the slow, expensive, judgement-bearing step.
This paper operationalizes that asymmetry end to end and shows that an autonomous
agent can carry the heavier half cheaply and reproducibly.

The work extends a documented lineage. The prior paper established a two-stage
VVUQ framework that placed verification over generation for a large language
model oncology code workflow, treating the assurance of generated code as the
deliverable rather than the code itself \cite{kawchak2026vvuq01}. The present
study specializes that generic framework in three ways. It narrows generic code
verification to ten humanoid-specific surgical gates, each tied to a named
failure mode of a two-handed standing robot. It binds every gate to a wired
corpus of robotic-surgery and service-robot safety standards resolved at runtime,
so the thresholds are not invented but inherited from published consensus
documents. And it exercises the whole apparatus over a concrete 60-second Whipple
procedure rather than an abstract code artifact. The study sits within a broader
program of autonomous physical AI oncology trial automation that includes a fully
automated trial sponsor \cite{kawchak2026sponsor}, a federated multi-site trial
unification layer \cite{kawchak2026federated}, and accelerated patient outcome
prediction \cite{kawchak2026prediction}, and it draws its pancreatic cancer
modelling heritage from a quantitative systems pharmacology in silico trial
\cite{kawchak2025qsppdac} and an end-to-end pancreatic ductal adenocarcinoma
digital twin proposal \cite{kawchak2025pdactwin}. The contribution here is to
make the assurance step itself the unit of automation.

The scope is deliberately frozen so the assurance argument stays clean. There is
one simulated patient, the borderline-resectable pancreatic ductal
adenocarcinoma case PAT-PDAC-0001, reused verbatim from the multi-arm baseline so
the single-humanoid run is directly comparable. There is one procedure, a classic
pancreaticoduodenectomy with pylorus preservation rendered as a 60-second
eight-phase timeline. There is one robot, the hypothetical 2030 Unitree
H2-Surgical 1.0 with two seven-degree-of-freedom arms and two
twenty-degree-of-freedom dexterous hands. And there is one deterministic design,
a 32-iteration sweep at seed 20260525. The timing is opportune. The United States
Food and Drug Administration has announced steps toward real-time clinical trials
\cite{fda2026realtime}, and its guidance on assessing the credibility of
computational modeling and simulation in medical device submissions
\cite{fda2023credibility} gives a concrete credibility framework that a
standards-anchored, autonomous assurance pipeline is well placed to support.
These references are made respectfully and without presumption; the argument is
simply that a fast, reproducible, standards-anchored assurance method is timely,
and that the United States is well positioned to lead the world in
surgical-humanoid VVUQ tied to external standards for patient safety, efficacy,
and trial speed.

This paper makes four contributions, each mapped to the Results section that
follows. First, it presents a ten-gate humanoid assurance decision surface and
demonstrates that its accept, block, and escalate paths all behave as specified
across five executed cases. Second, it reports a 32-iteration deterministic sweep
and a four-entrant comparison tournament that together show the assurance overlay
is decisive over any headline score. Third, it features two substantial generated
artifacts, a structured four-competitor results file and a thousand-row
non-repetitive humanoid sensor stream, and shows that producing them is cheap
while verifying them is the work. Fourth, it argues that binding every gate to
published external standards is what turns an inexpensive autonomous run into
defensible evidence. The remainder of the paper proceeds from Methods, which
describes the environment, the platform, the ten-gate design, the wired standards
corpus, the pipeline, and the verification floor, through Results, Discussion,
Limitations and Future Work, and Conclusions.
